\clearpage
\section{Sources and theorem boundaries}
The external source check distinguishes a stated theorem, a contextual reference and an unresolved comparison. The critical conjecture quantifiers, known unconditional bound, polynomial theorem, modified Szpiro equivalence, modularity, polylogarithm formulas and current IUT status have authenticated primary-source leaves in the evidence package.

The prime contact and determinant calculations are derived in this text. Pasten's arithmetic derivative work supplies an established base for that branch. Its small-derivative conjecture and its weak $abc$ equivalences are not silently identified with the stronger all-$\eps$ packet target stated here. The Wiles-defect and motivic references explain relevant established frameworks. They do not construct the project-specific persistence-to-Hecke or integral period transfer maps.

Otsubo and Yamazaki's motivic identities have precise coefficient-field and base-field hypotheses. Some Gauss-sum statements are over finite fields, while the multiplication relation extends to fields containing the needed roots of unity. None gives the missing integral radical-scale norm comparison in this work. Borel's rank input is used only for the rational weight motivation explicitly described in the path-packet section.

Bibliographic entries inherited solely as mathematical background are retained for that purpose. The exact source-check report states which leaves were freshly inspected. A failed search or a citation title is not treated as proof of a theorem or of historical novelty.

\begingroup
\begin{thebibliography}{99}
\raggedright
\bibitem{FaltingsIsogeny1986}
G. Faltings,
``Finiteness theorems for abelian varieties over number fields,''
in G. Cornell and J. H. Silverman (eds.), \emph{Arithmetic Geometry},
Springer, 1986, Chapter II, 9--26.
\href{https://doi.org/10.1007/978-1-4613-8655-1_2}{doi:10.1007/978-1-4613-8655-1\_2}.
Translation of the 1983 German original with its 1984 erratum.


\bibitem{SerreAbelian1989}
J.-P. Serre,
\emph{Abelian $\ell$-adic representations and elliptic curves},
with the collaboration of W. Kuyk and J. Labute,
Addison--Wesley, 1989, reprint of the 1968 original.
Chapter IV, Appendix A.1.2--A.1.3; Special Preface, p.~ix.


% Verified publication metadata and exact consulted versions, 2026-10-08.
% Bibitem fragments only: insert inside an existing thebibliography environment.
% Do not duplicate existing keys. See LITERATURE_CHECK.md for dependency scope.

\bibitem{StewartYu2001}
C. L. Stewart and K. R. Yu,
``On the $abc$ conjecture, II,''
\emph{Duke Mathematical Journal} 108, no.~1 (2001), 169--181.
\href{https://doi.org/10.1215/S0012-7094-01-10815-6}{doi:10.1215/S0012-7094-01-10815-6}.

\bibitem{Snyder2000}
N. Snyder,
``An alternate proof of Mason's theorem,''
\emph{Elemente der Mathematik} 55, no.~3 (2000), 93--94.
\href{https://doi.org/10.1007/S000170050074}{doi:10.1007/S000170050074}.

\bibitem{Oesterle1988}
J. Oesterl\'e,
``Nouvelles approches du `th\'eor\`eme' de Fermat,''
S\'eminaire Bourbaki, expos\'e 694,
\emph{Ast\'erisque} 161--162 (1988), 165--186.
\url{https://www.numdam.org/item/SB_1987-1988__30__165_0/}.

\bibitem{PastenDerivatives}
H. Pasten,
``Arithmetic derivatives through geometry of numbers,''
\emph{Canadian Mathematical Bulletin} 65, no.~4 (2022), 906--923.
\href{https://doi.org/10.4153/S0008439521000990}{doi:10.4153/S0008439521000990}.
Consulted author version: arXiv:2106.16165v3, 11 December 2021.

\bibitem{VojtaABC}
P. Vojta,
``A more general $ABC$ conjecture,''
\emph{International Mathematics Research Notices} 1998, no.~21, 1103--1116.
\href{https://doi.org/10.1155/S1073792898000658}{doi:10.1155/S1073792898000658}.

\bibitem{MochizukiIUTIII}
S. Mochizuki,
``Inter-universal Teichm\"uller theory III: Canonical splittings of the log-theta-lattice,''
\emph{Publications of the Research Institute for Mathematical Sciences}
57, no.~1/2 (2021), 403--626.
\href{https://doi.org/10.4171/PRIMS/57-1-3}{doi:10.4171/PRIMS/57-1-3}.
Consulted author PDF: May 2020.

\bibitem{MochizukiIUTIV}
S. Mochizuki,
``Inter-universal Teichm\"uller theory IV: Log-volume computations and set-theoretic foundations,''
\emph{Publications of the Research Institute for Mathematical Sciences}
57, no.~1/2 (2021), 627--723.
\href{https://doi.org/10.4171/PRIMS/57-1-4}{doi:10.4171/PRIMS/57-1-4}.
Consulted author PDF: April 2020.

\bibitem{ScholzeStix2018}
P. Scholze and J. Stix,
``Why $abc$ is still a conjecture,'' August 2018 manuscript.
\url{https://www.math.uni-bonn.de/people/scholze/WhyABCisStillaConjecture.pdf}.

\bibitem{MochizukiComments2018}
S. Mochizuki,
``Comments on the manuscript (2018-08 version) by Scholze--Stix concerning inter-universal Teichm\"uller theory (IUTch),'' September 2018.
\url{https://www.kurims.kyoto-u.ac.jp/~motizuki/Cmt2018-08.pdf}.

\bibitem{MochizukiReport2025}
S. Mochizuki,
``Report on the current situation surrounding inter-universal Teichm\"uller theory (IUT),'' October 2025.
\url{https://www.kurims.kyoto-u.ac.jp/~motizuki/IUT-report-2025-10.pdf}.

\bibitem{LANA2026}
LANA Project,
\emph{Project LANA interim report on IUT theory}, last updated 20 July 2026.
Author report: \url{https://github.com/katobungen/LANA_report_202607}.
Official announcement: \url{https://zen.ac.jp/news/zmcpostevent0717e}.

\bibitem{DLMFPolylog}
NIST Digital Library of Mathematical Functions,
\S25.12, ``Polylogarithms,'' equations 25.12.4, 25.12.6, 25.12.10, and 25.12.13.
and \S25.11, equation 25.11.14. Accessed 8 October 2026.
\url{https://dlmf.nist.gov/25.12}. \url{https://dlmf.nist.gov/25.11}.

\bibitem{Wiles}
A. Wiles,
``Modular elliptic curves and Fermat's Last Theorem,''
\emph{Annals of Mathematics} (2) 141, no.~3 (1995), 443--551.
\href{https://doi.org/10.2307/2118559}{doi:10.2307/2118559}.

\bibitem{TaylorWiles}
R. Taylor and A. Wiles,
``Ring-theoretic properties of certain Hecke algebras,''
\emph{Annals of Mathematics} (2) 141, no.~3 (1995), 553--572.
\href{https://doi.org/10.2307/2118560}{doi:10.2307/2118560}.

\bibitem{BCDT}
C. Breuil, B. Conrad, F. Diamond, and R. Taylor,
``On the modularity of elliptic curves over $\mathbb Q$: wild 3-adic exercises,''
\emph{Journal of the American Mathematical Society} 14, no.~4 (2001), 843--939.
\href{https://doi.org/10.1090/S0894-0347-01-00370-8}{doi:10.1090/S0894-0347-01-00370-8}.

\bibitem{Agashe}
A. Agashe,
``The modular number, congruence number, and multiplicity one,''
arXiv:0810.5176v2, 5 October 2025 (first version 2008).
\url{https://arxiv.org/abs/0810.5176v2}.

\bibitem{CesnaviciusManin}
K. \v{C}esnavi\v{c}ius, M. Neururer, and A. Saha,
``The Manin constant and the modular degree,''
\emph{Journal of the European Mathematical Society} 26, no.~2 (2024), 573--637.
\href{https://doi.org/10.4171/JEMS/1367}{doi:10.4171/JEMS/1367}.
Consulted arXiv:1911.09446v3, 2 November 2022.

\bibitem{BKM1}
G. B\"ockle, C. B. Khare, and J. Manning,
``Wiles defect for Hecke algebras that are not complete intersections,''
\emph{Compositio Mathematica} 157, no.~9 (2021), 2046--2088.
\href{https://doi.org/10.1112/S0010437X21007454}{doi:10.1112/S0010437X21007454}.
Consulted arXiv:1910.08507v3, 4 April 2021.

\bibitem{BKM2}
G. B\"ockle, C. B. Khare, and J. Manning,
``Wiles defect of Hecke algebras via local-global arguments,''
with an appendix by N. Fakhruddin and C. B. Khare,
\emph{Journal of the Institute of Mathematics of Jussieu} 23, no.~6 (2024), 2461--2541.
\href{https://doi.org/10.1017/S1474748024000021}{doi:10.1017/S1474748024000021}.
Consulted arXiv:2108.09729v2, 24 August 2021.

\bibitem{OtsuboYamazaki}
N. Otsubo and T. Yamazaki,
``Motivic Gauss and Jacobi sums,''
\emph{Annales de la Facult\'e des Sciences de Toulouse}, series 6,
35, no.~1 (2026), 95--135.
\href{https://doi.org/10.5802/afst.1842}{doi:10.5802/afst.1842}.
Consulted arXiv:2402.06072v3, 2 March 2025.

\bibitem{Goncharov}
A. B. Goncharov,
``Multiple polylogarithms and mixed Tate motives,''
arXiv:math/0103059v4, 15 May 2001.
\url{https://arxiv.org/abs/math/0103059v4}.

\bibitem{ZagierPolylog}
D. Zagier,
``Polylogarithms, Dedekind zeta functions, and the algebraic $K$-theory of fields,''
in \emph{Arithmetic Algebraic Geometry}, Progress in Mathematics 89,
Birkh\"auser, 1990, 391--430.
\href{https://doi.org/10.1007/978-1-4612-0457-2_19}{doi:10.1007/978-1-4612-0457-2\_19}.

\bibitem{Borel}
A. Borel,
``Stable real cohomology of arithmetic groups,''
\emph{Annales scientifiques de l'\'Ecole Normale Sup\'erieure}, series 4,
7, no.~2 (1974), 235--272.
\href{https://doi.org/10.24033/asens.1269}{doi:10.24033/asens.1269}.

\bibitem{SUnit1988}
J.-H. Evertse, K. Gy\H{o}ry, C. L. Stewart, and R. Tijdeman,
``On $S$-unit equations in two unknowns,''
\emph{Inventiones Mathematicae} 92 (1988), 461--477.
\href{https://doi.org/10.1007/BF01393743}{doi:10.1007/BF01393743}.

\bibitem{AokiFermat}
N. Aoki,
``Simple factors of the Jacobian of a Fermat curve and the Picard number of a product of Fermat curves,''
\emph{American Journal of Mathematics} 113 (1991), 779--833.
\href{https://doi.org/10.2307/2374786}{doi:10.2307/2374786}.

\bibitem{ColemanMcCallum}
R. Coleman and W. McCallum,
``Stable reduction of Fermat curves and Jacobi sum Hecke characters,''
\emph{Journal f\"ur die reine und angewandte Mathematik} 385 (1988), 41--101.
\href{https://doi.org/10.1515/crll.1988.385.41}{doi:10.1515/crll.1988.385.41}.

\end{thebibliography}
\endgroup
